\section{Related work}
\label{sec:related}

\paragraph{Agents that act, observe, and try again.}
ReAct \citep{yao2023react} established the template that dominates practice:
interleave a written action with the environment's observation and let the model
condition on the growing transcript. Reflexion \citep{shinn2023reflexion},
Self-Refine \citep{madaan2023selfrefine} and self-debugging
\citep{chen2024selfdebug} all build on it, adding an explicit critique or a test
result to the same append-only transcript. MINT \citep{wang2024mint} evaluates
multi-turn interaction with tools and language feedback directly and finds
feedback helpful in aggregate. The shared premise is that appending the failure
is informative; what none of these measure is the counterfactual we study here,
namely what the appended \emph{action string} does on its own.

\paragraph{Doubts about self-correction.}
A parallel line argues that models cannot reliably correct themselves without an
external signal \citep{huang2024selfcorrect,valmeekam2023critiquing}. Our result
is not a restatement of that. Those papers ask whether a model can \emph{judge}
its own output; we hold the judgement fixed, since the environment executes the
action and the verdict is ground truth, and ask whether the model's next-token
distribution moves in the direction the verdict implies. The failure we document
survives perfect feedback.

\paragraph{Context contamination.}
The closest prior work is \citet{yang2026contamination}, who formalises the
observation that a retry in a context still containing the failed attempt has a
higher error rate than a clean one. The Context-Contaminated Restart Model gives
closed-form results for success under a budget and fits SWE-bench Verified data
with a cascade ratio $\epsilon_1/\epsilon_0 = 7.1$, showing that assuming
independent retries overestimates pass@3 by 17.4 points. That work establishes
the phenomenon at the level of task statistics and concludes that the remedy is
to \emph{clear the context before retrying}. Our contribution is orthogonal and,
we think, complementary in a useful way: we ask which \emph{part} of the context
does the contaminating, measure it at the token level, and find that it is the
failed action's surface form rather than the error message. That matters
practically, because clearing the context discards the diagnosis along with the
hazard. In our terms, clean restart is the \emph{drop} harness, and it is a
baseline we evaluate rather than an endpoint.

\paragraph{Agents that get stuck.}
Repetition in agent trajectories is widely reported and usually treated as an
engineering defect. \citet{hou2026infiniteloops} detect unbounded feedback paths
by static analysis of agent \emph{source code} across 6{,}549 repositories,
which is complementary to our question: their loops are properties of the
program, ours of the policy. \citet{wang2025hellorhighwater} show that agents
struggle to formulate alternative plans after an external failure even when the task is guaranteed to remain solvable, which is a behavioural counterpart
to what we measure distributionally. Most striking is \citet{zenkri2026fidelity}, who find
on a mechanical puzzle that \emph{degrading} an embodied agent's observations
improves success up to $2.85\times$, and trace the gain to fewer repetitive
action loops. That is independent evidence that the observation channel can be
net harmful; we supply the token-level mechanism and a fix that does not require
corrupting the observations.

\paragraph{Repetition and copying in language models.}
\citet{xu2022breaktheloop} showed that sentence repetition is
self-reinforcing, in that the more often a sentence appears in the context the
more likely the model is to produce it again, and proposed a training-time
penalty.
Degeneration \citep{holtzman2020degeneration}, unlikelihood training
\citep{welleck2020unlikelihood} and repetition penalties \citep{keskar2019ctrl}
address the same tendency in open-ended generation. Mechanistically, induction
heads \citep{olsson2022induction} implement prefix matching and copying, copy
suppression \citep{mcdougall2024copysuppression} implements a counterweight, and
\citet{repetitionneurons2025} isolate components whose over-dominance yields the
``repetition curse''. All of this concerns generation \emph{without} a competing
signal. The agent loop is the interesting case precisely because a competing
signal is present by design, and the question is which one wins.

\paragraph{Negative instructions.}
Naming something in a prompt raises its probability even under an instruction
not to produce it, a pattern documented for topic avoidance
\citep{castricato2024pinkelephants,nopinkelephant2024} and for negation more
broadly \citep{negation2025pinkelephant}. Our $\copyterm$ is the same family of
effect measured inside a trajectory rather than a prompt, and
Section~\ref{sec:agent} shows that the corresponding remedy, adding a
``do not repeat'' instruction, behaves exactly as that literature predicts.

\paragraph{Small models as agents.}
\citet{belcak2025slmagents} argue that most agentic sub-tasks are within reach of
small models, making the harness the deciding factor.
\citet{toolschemas2025} adapt tool schemas rather than models, and
\citet{constrainttax2026} document that structured-output constraints can
suppress tool calling outright. Constrained decoding is well developed as a
\emph{format} device \citep{willard2023outlines}, with known costs to free-form
reasoning \citep{tam2024speakfreely}. We use it differently: not to enforce a
grammar, but to make one specific previously-failed string unreachable. To our
knowledge the ban-what-already-failed use of constrained decoding has not been
evaluated as an agent-reliability intervention.

\paragraph{Spending more compute at test time.}
A large literature buys accuracy with repeated attempts
\citep{brown2024monkeys,snell2024testtime}, subject to the quality of whatever
selects among them \citep{stroebl2024inferenceflaws}. Our result is a constraint
on that arithmetic in the sequential case. Attempts within a single trajectory
are not independent draws: each one is conditioned on a context containing its
predecessors, and we measure that conditioning pushing toward repetition. The
same point is made from the outcome side by \citet{yang2026contamination}, whose
fit implies that assuming independence overestimates pass@3 substantially. If
extra attempts are to pay, the harness has to make them differ.

\paragraph{Context management.}
\citet{liu2024lostmiddle} showed that position within a long context matters;
recent agent work learns to curate context, keeping compact state that includes
records of ineffective attempts \citep{adacom2026context}. Those systems edit
context to save tokens or to keep salient constraints in view. Our analysis
implies a different reason to edit it, namely that the verbatim failed action
is itself the hazard, and predicts which edit helps, which we test directly.

\paragraph{Positioning.}
Table~\ref{tab:related} summarises the difference. Prior work establishes that
(i) repetition is self-reinforcing without feedback, (ii) agents empirically get
stuck, and (iii) negative instructions are weak. What is missing, and what we
provide, is the controlled measurement of whether execution feedback moves the
action distribution in the right direction at all, a decomposition of that
quantity into the parts a harness can and cannot change, and an intervention
derived from the decomposition rather than from intuition.

\begin{table*}[t]
\centering
\footnotesize
\setlength{\tabcolsep}{4pt}
\begin{tabular}{l c c c}
\toprule
 & \makecell{explicit\\failure signal} & \makecell{separates form\\from meaning} & \makecell{harness\\intervention} \\
\midrule
Repetition / degeneration                  & no  & no  & decoding \\
\quad \citep{xu2022breaktheloop,holtzman2020degeneration} & & & \\
Negative instructions                      & n/a & no  & prompt \\
\quad \citep{castricato2024pinkelephants}  & & & \\
Self-correction critiques                  & yes & no  & prompt \\
\quad \citep{huang2024selfcorrect,shinn2023reflexion} & & & \\
Agent loop detection                       & yes & no  & static analysis \\
\quad \citep{hou2026infiniteloops}         & & & \\
Observation-fidelity probe                 & yes & no  & degrade inputs \\
\quad \citep{zenkri2026fidelity}           & & & \\
Context contamination (CCRM)               & yes & no  & clear the context \\
\quad \citep{yang2026contamination}        & & & \\
\midrule
This work                                  & yes & \textbf{yes} & context + decoder \\
\bottomrule
\end{tabular}
\caption{Where this paper sits. The distinguishing column is the middle one:
prior work measures the \emph{net} effect of a failure record, if it measures it
at all, whereas the intervention that works depends on which of its two
components dominates.}
\label{tab:related}
\end{table*}

